%%
%% paper77_conway_knot_sliceness_en.tex
%% Autor: Michael Fuhrmann
%% Projekt: Specialist Mathematics Project
%% Thema: Conway's Knot and the Sliceness Problem in Knot Theory
%% Sprache: Englisch
%% Erstellt: 2026-03-12
%% Letzte Änderung: 2026-03-12
%% Build: 164
%%
%% Beschreibung:
%%   Dieses Paper behandelt das Conway-Knoten-Problem. Der Conway-Knoten
%%   (11 Kreuzungen) war das einzige Prime-Knoten mit ≤ 12 Kreuzungen,
%%   dessen glatte Schnittigkeit unbekannt war. Lisa Piccirillo bewies 2020,
%%   dass der Conway-Knoten NICHT glatt scheibig ist. Die topologische
%%   Schnittigkeit bleibt ein offenes Problem.
%%

\documentclass[12pt,a4paper]{amsart}

%% Pakete für Mathematik, Layout, Encoding und Links
\usepackage{amsmath,amssymb,amsthm,mathtools,enumitem,booktabs,array,hyperref}
\usepackage[utf8]{inputenc}
\usepackage[T1]{fontenc}
\usepackage{geometry}
\geometry{margin=2.5cm}

%% Theorem-Umgebungen
\newtheorem{theorem}{Theorem}[section]
\newtheorem{lemma}[theorem]{Lemma}
\newtheorem{corollary}[theorem]{Corollary}
\newtheorem{proposition}[theorem]{Proposition}
\theoremstyle{definition}
\newtheorem{definition}[theorem]{Definition}
\newtheorem{example}[theorem]{Example}
\theoremstyle{remark}
\newtheorem{remark}[theorem]{Remark}
\newtheorem{conjecture}[theorem]{Conjecture}

%% Mathematische Operatoren
\DeclareMathOperator{\Arf}{Arf}
\DeclareMathOperator{\sign}{sign}
\DeclareMathOperator{\rank}{rank}
\DeclareMathOperator{\HF}{HF}
\DeclareMathOperator{\genus}{genus}
\newcommand{\Z}{\mathbb{Z}}
\newcommand{\Q}{\mathbb{Q}}
\newcommand{\R}{\mathbb{R}}
\newcommand{\C}{\mathbb{C}}

%% Hyperref-Konfiguration
\hypersetup{
    colorlinks=true,
    linkcolor=blue,
    citecolor=blue,
    urlcolor=blue,
    pdftitle={Conway's Knot and the Sliceness Problem in Knot Theory},
    pdfauthor={Michael Fuhrmann}
}

%% -------------------------------------------------------------------------
\begin{document}
%% Abstract VOR \maketitle
%% -------------------------------------------------------------------------
\begin{abstract}
The Conway knot $C$, an 11-crossing knot discovered by John Horton Conway,
is the mutant of the Kinoshita--Terasaka knot. For decades it was the only
prime knot with 12 or fewer crossings whose smooth concordance class was
undetermined. In 2020, Lisa Piccirillo~\cite{Piccirillo2020} resolved this
problem by proving that $C$ is \emph{not} smoothly slice, using a remarkable
construction: she found a knot $K$ with the same zero-surgery $3$-manifold
as $C$ but with non-vanishing Rasmussen $s$-invariant, then applied Freedman's
theorem on topologically slice knots with vanishing Arf invariant to produce
a contradiction.

We present the concordance group $\mathcal{C}$, the distinction between
topological and smooth sliceness, the relevant concordance invariants
(Arf invariant, Rasmussen $s$-invariant, Heegaard Floer $\tau$-invariant),
Piccirillo's proof in detail, and the remaining open questions — in particular,
whether the Conway knot is topologically slice.
\end{abstract}

\title{Conway's Knot and the Sliceness Problem in Knot Theory}
\author{Michael Fuhrmann}
\address{Specialist Mathematics Project, Build 165}
\email{specialist-maths@project.local}
\date{2026-03-12}

\maketitle

%% -------------------------------------------------------------------------
\section{Introduction}
%% -------------------------------------------------------------------------

Knot concordance is one of the central organising themes of low-dimensional
topology. Two knots $K_0, K_1 \subset S^3$ are \emph{concordant} if there
exists a smoothly embedded annulus $A \cong S^1 \times [0,1]$ in
$S^3 \times [0,1]$ with $\partial A = K_1 \times \{1\} \cup (-K_0) \times \{0\}$.
Concordance classes form a group $\mathcal{C}$ under connected sum, the
\emph{concordance group}.

A knot is \emph{slice} if it is concordant to the unknot, equivalently if
it bounds a smoothly embedded disc in $B^4$. The distinction between
\emph{topologically slice} (disk exists as a topological locally flat embedding)
and \emph{smoothly slice} (smooth embedding) is one of the most profound
manifestations of the difference between the topological and smooth
categories in dimension 4.

The \emph{Conway knot} $C$ (Figure~1) is the unique 11-crossing prime knot
denoted $11n34$ in the Hoste--Thistlethwaite table. It is the mutant
of the Kinoshita--Terasaka knot $KT$ ($11n42$), obtained by cutting along
a Conway sphere and regluing with a $180°$ rotation. The Conway knot
satisfies all known necessary conditions for being smoothly slice: its
Arf invariant vanishes, its Alexander polynomial satisfies the Fox--Milnor
factorisation, and it is algebraically slice. Yet for decades, its smooth
concordance class remained undetermined.

In a landmark paper~\cite{Piccirillo2020}, Lisa Piccirillo (2020) proved:

\begin{theorem}[Piccirillo 2020~\cite{Piccirillo2020}]
\label{thm:Piccirillo}
The Conway knot is not smoothly slice.
\end{theorem}

Piccirillo's proof is a tour de force of 4-manifold topology, using
Freedman's theorem, Kirby calculus, and the Rasmussen $s$-invariant.

%% -------------------------------------------------------------------------
\section{Concordance and Sliceness}
%% -------------------------------------------------------------------------

\subsection{The Concordance Group}

\begin{definition}[Knot Concordance]
Two oriented knots $K_0, K_1 \subset S^3$ are \emph{smoothly concordant},
written $K_0 \sim K_1$, if there exists a smooth proper embedding
$f \colon S^1 \times [0,1] \hookrightarrow S^3 \times [0,1]$
with $f(S^1 \times \{i\}) = K_i$ for $i = 0, 1$.
\end{definition}

\begin{theorem}[Fox--Milnor~\cite{FoxMilnor1966}]
The set of concordance classes of oriented knots in $S^3$ forms an
abelian group $\mathcal{C}$ under connected sum $\#$, the
\emph{smooth concordance group}. The identity is the class of the unknot
(= the class of slice knots). The inverse of $[K]$ is $[-K]$
(the mirror image with reversed orientation).
\end{theorem}

\subsection{Topological vs.\ Smooth Sliceness}

\begin{definition}[Topologically and Smoothly Slice]
A knot $K \subset S^3$ is:
\begin{enumerate}[label=\roman*)]
\item \emph{smoothly slice} if it bounds a smoothly embedded $2$-disc
      $D^2 \hookrightarrow B^4$;
\item \emph{topologically slice} if it bounds a locally flat topologically
      embedded $2$-disc $D^2 \hookrightarrow B^4$.
\end{enumerate}
Every smoothly slice knot is topologically slice. The converse fails:
there exist topologically but not smoothly slice knots.
\end{definition}

\begin{example}
Whitehead doubles of non-slice knots provide the standard examples of
topologically but not smoothly slice knots: for a non-slice knot $J$,
the (untwisted) Whitehead double $\mathrm{Wh}(J)$ is topologically slice
(since $\Delta_{\mathrm{Wh}(J)}(t) = 1$, Freedman's theorem applies) but
not smoothly slice, as detected by the Rasmussen $s$-invariant or
Heegaard Floer $\tau$-invariant~\cite{Hedden2007}.
Note: torus knots $T(2, 2k+1)$ are \emph{not} topologically slice
(their signatures are non-zero and $\Delta_{T(2,2k+1)}(t) \neq 1$).
\end{example}

%% -------------------------------------------------------------------------
\section{The Conway Knot and its Mutant}
%% -------------------------------------------------------------------------

\subsection{Conway Spheres and Mutation}

\begin{definition}[Conway Sphere and Mutation]
A \emph{Conway sphere} in $S^3 \setminus K$ is an embedded $2$-sphere
meeting $K$ in exactly 4 points, dividing $K$ into two $2$-tangles.
The \emph{mutant} $K^m$ of $K$ is obtained by cutting along the Conway
sphere and regluing by the involution that rotates the sphere $180°$
around a horizontal or vertical axis.
\end{definition}

The Conway knot $C$ and the Kinoshita--Terasaka knot $KT$ are related by
mutation: they share the same Conway sphere, and $C$ is obtained from $KT$
by a mutation. Several concordance invariants are \emph{mutation-invariant}:

\begin{theorem}
The following concordance invariants are preserved under mutation:
\begin{enumerate}[label=\roman*)]
\item Alexander polynomial $\Delta_K(t)$;
\item Signature $\sigma(K)$;
\item Arf invariant $\Arf(K)$;
\item Determinant $\det(K) = |\Delta_K(-1)|$.
\end{enumerate}
\end{theorem}

Since the Kinoshita--Terasaka knot $KT$ is slice (it bounds a ribbon disc
in $B^4$~\cite{KinoshitaTerasaka1957}), all these invariants vanish for $KT$,
and hence for its mutant $C$ as well.

\begin{proposition}
The Conway knot $C$ is algebraically slice: there exist generators
$\{a_i, b_i\}$ of $H_1(\Sigma_2(C))$ with $\mathrm{lk}(a_i, b_j^+) = 0$,
where $\Sigma_2(C)$ denotes the double branched cover.
\end{proposition}

\begin{proof}
The Seifert matrix of $C$ has a metaboliser (a rank-$g$ subgroup of
$H_1(F; \Z)$ on which the Seifert pairing vanishes), which is the algebraic
sliceness condition of Levine~\cite{Levine1969}.
\end{proof}

%% -------------------------------------------------------------------------
\section{Concordance Invariants}
%% -------------------------------------------------------------------------

\subsection{The Arf Invariant}

\begin{definition}[Arf Invariant]
The \emph{Arf invariant} of $K$ is defined as $\Arf(K) \in \Z/2\Z$.
It can be computed as $\Arf(K) = \Delta_K(-1) \pmod 2 / 2$ for odd
$\Delta_K(-1)$, or equivalently via the Goeritz matrix modulo 8.
\end{definition}

\begin{theorem}[Robertello 1965~\cite{Robertello1965}]
If $K$ is slice, then $\Arf(K) = 0$.
\end{theorem}

For the Conway knot: $\Arf(C) = 0$, so this obstruction vanishes.

\subsection{The Alexander Polynomial and Fox--Milnor Condition}

\begin{theorem}[Fox--Milnor~\cite{FoxMilnor1966}]
If $K$ is slice, then $\Delta_K(t) = f(t) \cdot f(t^{-1})$ for some
$f(t) \in \Z[t, t^{-1}]$.
\end{theorem}

The Alexander polynomial of the Conway knot satisfies this condition
(as does that of $KT$, since they are the same polynomial).

\subsection{The Rasmussen \texorpdfstring{$s$}{s}-Invariant}
\label{sec:s-invariant}

\begin{definition}[$s$-Invariant~\cite{Rasmussen2010}]
The \emph{Rasmussen $s$-invariant} is a knot concordance invariant
$s \colon \mathcal{K} \to 2\Z$ derived from Khovanov homology.
It is a homomorphism on $\mathcal{C}$ and satisfies:
\begin{enumerate}[label=\roman*)]
\item $s(K) = 0$ if $K$ is slice;
\item $|s(K)| \leq 2 g_4(K)$ where $g_4$ is the $4$-ball genus;
\item $s(T(p,q)) = (p-1)(q-1)$ for torus knots.
\end{enumerate}
\end{definition}

\begin{theorem}[Rasmussen 2010~\cite{Rasmussen2010}]
$s(K)/2$ gives a lower bound on the smooth $4$-ball genus $g_4(K)$:
\[
    |s(K)|/2 \leq g_4(K).
\]
In particular, $s(K) = 0$ is necessary for $K$ to be smoothly slice.
\end{theorem}

The $s$-invariant is \emph{not} mutation-invariant in general:
Piccirillo's key insight was that the Conway knot has a ``companion knot''
with non-zero $s$-invariant.

\subsection{The Heegaard Floer \texorpdfstring{$\tau$}{tau}-Invariant}

\begin{definition}[$\tau$-Invariant~\cite{OzsvathSzabo2003}]
The Ozsváth--Szabó \emph{$\tau$-invariant} is a concordance homomorphism
$\tau \colon \mathcal{C} \to \Z$ defined via the knot Floer homology
$\widehat{HFK}(K)$. It satisfies:
\begin{enumerate}[label=\roman*)]
\item $\tau(K) = 0$ for slice knots;
\item $|\tau(K)| \leq g_4(K)$;
\item $\tau(T(p,q)) = (p-1)(q-1)/2$.
\end{enumerate}
\end{definition}

Both $\tau$ and $s$ agree on ``standard'' knots, but they are conjectured
to differ on certain knots (and indeed do differ; see~\cite{Hedden2009}).

%% -------------------------------------------------------------------------
\section{Piccirillo's Proof}
%% -------------------------------------------------------------------------

We now present the main ideas of Piccirillo's proof~\cite{Piccirillo2020}
that the Conway knot is not smoothly slice.

\subsection{The Key Construction}

\begin{definition}[Zero Surgery]
The \emph{zero surgery} $S^3_0(K)$ is the 3-manifold obtained from $S^3$
by Dehn surgery on $K$ with framing coefficient 0.
\end{definition}

\begin{lemma}[Piccirillo~\cite{Piccirillo2020}]
\label{lem:Piccirillo}
Let $K'$ be the knot obtained from $C$ by a specific modification
via the Kirby diagram of an exotic $B^4$ constructed using the Conway
knot. Then:
\begin{enumerate}[label=\roman*)]
\item $S^3_0(K') \cong S^3_0(C)$ (as oriented 3-manifolds).
\item $s(K') \neq 0$.
\end{enumerate}
\end{lemma}

The construction of $K'$ uses the following strategy:
Start with the trace $X_C = B^4 \cup_C h^2$ (the 4-manifold with boundary
$S^3_0(C)$ obtained by attaching a 2-handle along $C$ with framing 0).
Piccirillo finds a different $2$-handle $h'^2$ attached along $K'$
with the same framing such that $X_{K'} \cong X_C$ (diffeomorphic
4-manifolds). This is a purely Kirby-calculus argument involving
a sequence of handle slides and cancellations.

\subsection{Freedman's Theorem}

\begin{theorem}[Freedman 1982~\cite{Freedman1982}]
Any knot $K \subset S^3$ whose Alexander polynomial is trivial
($\Delta_K(t) = 1$) is topologically slice.
\end{theorem}

\begin{remark}
Note that a vanishing Arf invariant alone does \emph{not} suffice for
topological sliceness. Freedman's theorem requires the stronger condition
$\Delta_K(t) = 1$ (equivalently, $\pi_1(S^3 \setminus K) \cong \Z$).
\end{remark}

The stronger version relevant here is:

\begin{theorem}[Freedman~\cite{Freedman1982}]
If $\pi_1(S^3 \setminus K) \cong \Z$ (i.e.\ $K$ is a homotopy-trivial
knot, which holds iff $\Delta_K = 1$) then $K$ is topologically slice.
\end{theorem}

\subsection{The Proof of Theorem~\ref{thm:Piccirillo}}

\begin{proof}[Proof that $C$ is not smoothly slice]
\textbf{Step 1:} Piccirillo constructs the knot $K'$ of Lemma~\ref{lem:Piccirillo}
explicitly via Kirby calculus on the Kirby diagram of $X_C$.

\textbf{Step 2:} She computes $s(K') \neq 0$ by computing the Khovanov
homology of $K'$ using a computer program (KnotTheory package in
\textit{Mathematica}). Specifically, $s(K') = -2$, which in particular shows
$K'$ is not smoothly slice by property (i) of the $s$-invariant.

\textbf{Step 3:} Now suppose for contradiction that $C$ were smoothly slice,
bounding a smooth disc $D \subset B^4$.

\textbf{Step 4:} Since $S^3_0(K') \cong S^3_0(C)$, there is a diffeomorphism
$\phi \colon X_{K'} \to X_C$ of the 4-manifold traces.

\textbf{Step 5:} Using the smooth disc $D$ for $C$, one constructs a smooth
disc for $K'$ by composing $\phi$ with the cobordism induced by $D$.
Concretely: the trace $X_C$ embeds into $B^4$ (since $C$ is assumed slice
and the slice disc provides a nullhomologous sphere); pulling back via
$\phi$ gives a smooth disc for $K'$ in $B^4$.

\textbf{Step 6:} But $K'$ is not smoothly slice ($s(K') = -2 \neq 0$).
This is a contradiction.

Therefore $C$ is not smoothly slice. \qed
\end{proof}

\begin{remark}
The proof uses no deep gauge theory or Floer homology — only:
(1) a Kirby calculus computation (constructing $K'$),
(2) a combinatorial computation ($s(K') \neq 0$),
(3) Freedman's topological theorem.
This is what makes the proof so elegant and unexpected.
\end{remark}

%% -------------------------------------------------------------------------
\section{The 4-Ball Genus and Ribbon Knots}
%% -------------------------------------------------------------------------

\begin{definition}[4-Ball Genus and Ribbon Genus]
The \emph{smooth $4$-ball genus} of $K$ is
\[
    g_4(K) = \min\{ g(F) \mid F \subset B^4 \text{ smooth, } \partial F = K \},
\]
where the minimum is over all compact smooth orientable surfaces $F$ in $B^4$.
$K$ is slice iff $g_4(K) = 0$.

$K$ is \emph{ribbon} if it bounds a ribbon disc (an immersed disc in $S^3$
with only ribbon singularities). Every ribbon knot is smoothly slice.
\end{definition}

\begin{conjecture}[Ribbon--Slice Conjecture]
\label{conj:ribbon-slice}
Every smoothly slice knot is ribbon.
\end{conjecture}

The Kinoshita--Terasaka knot $KT$ is known to be ribbon (and hence slice);
the Conway knot $C$ is not slice (Piccirillo), hence not ribbon.

\begin{theorem}[Piccirillo 2020]
$g_4(C) \geq 1$.
\end{theorem}

The exact value of $g_4(C)$ is not known; it is suspected to equal 1, but
this requires either exhibiting an explicit genus-1 surface or proving that
$g_4(C) \leq 1$.

%% -------------------------------------------------------------------------
\section{Topological Sliceness of the Conway Knot}
%% -------------------------------------------------------------------------

The question of \emph{topological} sliceness of the Conway knot is
substantially more subtle.

\begin{proposition}
The Conway knot satisfies the classical necessary conditions for topological
sliceness:
\begin{enumerate}[label=\roman*)]
\item $\Arf(C) = 0$.
\item $\Delta_C(t) = f(t) \cdot f(t^{-1})$ (Fox--Milnor condition).
\item $\sigma_\omega(C) = 0$ for all $\omega \in S^1$ (Levine-Tristram signatures vanish).
\item $C$ is algebraically slice (Witt class in Witt group of $\Q(t)$ is zero).
\end{enumerate}
\end{proposition}

\begin{conjecture}[Topological Sliceness of Conway Knot]
\label{conj:Conway-top-slice}
The Conway knot $C$ is topologically slice.
\end{conjecture}

This conjecture remains open as of 2026. The main obstacle to proving it
is the difficulty of applying Freedman's disk theorem in this setting:
Freedman's theorem requires the fundamental group of the knot complement to
be ``good'' in the 4-manifold sense, which holds for knots with
$\pi_1 \cong \Z$ (i.e.\ trivial Alexander polynomial), but the Conway knot
has non-trivial Alexander polynomial.

\begin{conjecture}[Topological Concordance vs.\ Smooth Concordance]
\label{conj:top-vs-smooth}
The Conway knot has a non-trivial image in the topological concordance group
$\mathcal{C}^{\mathrm{top}}$, i.e.\ it is not topologically concordant to
the unknot.
\end{conjecture}

If Conjecture~\ref{conj:Conway-top-slice} is false (i.e.\ $C$ is not
topologically slice), this would provide a new example of a knot that
``looks'' slice algebraically but is not even topologically slice — a rare
phenomenon.

%% -------------------------------------------------------------------------
\section{Ribbon Concordance}
%% -------------------------------------------------------------------------

\begin{definition}[Ribbon Concordance]
A concordance $A \colon K_0 \to K_1$ is \emph{ribbon} if the projection
$A \subset S^3 \times [0,1] \to [0,1]$ has only saddle-type critical points
(no local maxima), i.e.\ the concordance can be built by adding only
``saddle moves.''
\end{definition}

\begin{theorem}[Gordon 1981~\cite{Gordon1981}]
Ribbon concordance defines a partial order on the set of knots.
\end{theorem}

\begin{conjecture}[Gordon 1981~\cite{Gordon1981}]
\label{conj:Gordon}
Ribbon concordance is a \emph{total} order on prime knots: if $K_0$ and
$K_1$ are ribbon concordant in both directions, then $K_0 = K_1$.
\end{conjecture}

Recent progress by Agol~\cite{Agol2022} proved that ribbon concordance
induces an order on homotopy classes of 3-manifolds.

%% -------------------------------------------------------------------------
\section{Further Invariants and Methods}
%% -------------------------------------------------------------------------

\subsection{Heegaard Floer Homology Obstructions}

Knot Floer homology $\widehat{HFK}(K)$, developed by Ozsváth--Szabó
and independently Rasmussen, provides the $\tau$-invariant and the
$d$-invariants of surgery manifolds. For the Conway knot:

\begin{proposition}
$\tau(C) = 0$.
\end{proposition}

This was verified computationally and follows from the mutation-invariance
of $\tau$ (a theorem due to Ozsváth--Szabó~\cite{OzsvathSzabo2004}).
Since $\tau(KT) = 0$ (as $KT$ is slice), we have $\tau(C) = 0$.

\begin{remark}
The $s$-invariant is also conjecturally mutation-invariant, which would
give $s(C) = s(KT) = 0$ automatically. This is consistent with Piccirillo's
result (she only uses $s(K') \neq 0$ for her auxiliary knot $K'$, not $s(C)$).
\end{remark}

\subsection{Satellite Concordance Invariants}

For a pattern $P$ in a solid torus and companion knot $K$, the satellite
$P(K)$ satisfies:
\[
    \tau(P(K)) = w(P) \cdot \tau(K) + \tau(P(U)),
\]
where $w(P)$ is the winding number and $U$ is the unknot. Similar formulas
exist for $s$ under certain hypotheses on $P$ (Hedden~\cite{Hedden2009}).

These satellite formulas allow the detection of concordance information
from simpler invariants.

%% -------------------------------------------------------------------------
\section{Open Problems}
%% -------------------------------------------------------------------------

\begin{conjecture}[Topological Sliceness of $C$, Conjecture~\ref{conj:Conway-top-slice}]
The Conway knot is topologically slice.
\end{conjecture}

\begin{conjecture}[Ribbon--Slice, Conjecture~\ref{conj:ribbon-slice}]
Every smoothly slice knot is ribbon.
\end{conjecture}

\begin{conjecture}[Gordon's Conjecture, Conjecture~\ref{conj:Gordon}]
Ribbon concordance is a partial order (injective up to knot equivalence).
\end{conjecture}

\begin{conjecture}[Smooth 4-Genus of $C$]
$g_4(C) = 1$.
\end{conjecture}

\begin{conjecture}[Mutation and $s$-Invariant]
The Rasmussen $s$-invariant is preserved under mutation.
\end{conjecture}

\begin{conjecture}[Smooth Concordance Group]
The concordance group $\mathcal{C}$ contains a $\Z^\infty$ summand and
$\Z/2\Z$ torsion, but its full structure (in particular, whether it contains
any odd-order torsion) is unknown.
\end{conjecture}

%% -------------------------------------------------------------------------
\section{Conclusion}
%% -------------------------------------------------------------------------

The story of the Conway knot illustrates how the gap between the smooth and
topological categories in dimension 4 can be probed through concrete knot
theory. The resolution of its smooth sliceness by Piccirillo in 2020 is a
triumph of clever Kirby calculus combined with modern concordance invariants.
The central open question — whether $C$ is topologically slice — requires
entirely different methods, since all classical algebraic obstructions vanish,
and Freedman's geometric techniques do not apply directly. New ideas in
4-manifold topology will be needed to resolve this problem.

\begin{thebibliography}{99}

\bibitem{Piccirillo2020}
L.~Piccirillo,
\textit{The Conway knot is not slice},
Ann.\ of Math.\ \textbf{191} (2020), 581--591.

\bibitem{Conway1970}
J.~H.~Conway,
\textit{An enumeration of knots and links, and some of their algebraic
properties},
in: J.~Leech (ed.), \textit{Computational Problems in Abstract Algebra},
Pergamon Press, Oxford, 1970, 329--358.

\bibitem{KinoshitaTerasaka1957}
S.~Kinoshita, H.~Terasaka,
\textit{On unions of knots},
Osaka Math.\ J.\ \textbf{9} (1957), 131--153.

\bibitem{FoxMilnor1966}
R.~H.~Fox, J.~W.~Milnor,
\textit{Singularities of 2-spheres in 4-space and cobordism of knots},
Osaka J.\ Math.\ \textbf{3} (1966), 257--267.

\bibitem{Rasmussen2010}
J.~Rasmussen,
\textit{Khovanov homology and the slice genus},
Invent.\ Math.\ \textbf{182} (2010), 419--447.

\bibitem{OzsvathSzabo2003}
P.~Ozsváth, Z.~Szabó,
\textit{Knot Floer homology and the four-ball genus},
Geom.\ Topol.\ \textbf{7} (2003), 615--639.

\bibitem{OzsvathSzabo2004}
P.~Ozsváth, Z.~Szabó,
\textit{Knot Floer homology and integer surgeries},
Algebr.\ Geom.\ Topol.\ \textbf{8} (2008), 101--153.

\bibitem{Freedman1982}
M.~H.~Freedman,
\textit{The topology of four-dimensional manifolds},
J.\ Differential Geom.\ \textbf{17} (1982), 357--453.

\bibitem{KronheimerMrowka1993}
P.~B.~Kronheimer, T.~S.~Mrowka,
\textit{Gauge theory for embedded surfaces, I},
Topology \textbf{32} (1993), 773--826.

\bibitem{Levine1969}
J.~P.~Levine,
\textit{Knot cobordism groups in codimension two},
Comment.\ Math.\ Helv.\ \textbf{44} (1969), 229--244.

\bibitem{Robertello1965}
R.~A.~Robertello,
\textit{An invariant of knot cobordism},
Comm.\ Pure Appl.\ Math.\ \textbf{18} (1965), 543--555.

\bibitem{Gordon1981}
C.~McA.~Gordon,
\textit{Some aspects of classical knot theory},
in: \textit{Knot Theory}, Lecture Notes in Math.\ \textbf{685},
Springer, Berlin, 1978, 1--60.

\bibitem{Agol2022}
I.~Agol,
\textit{Ribbon concordance of knots is a partial order},
Ann.\ of Math.\ \textbf{196} (2022), 675--698.

\bibitem{Hedden2009}
M.~Hedden,
\textit{Notions of positivity and the Ozsváth-Szabó concordance invariant},
J.\ Knot Theory Ramifications \textbf{19} (2010), 617--629.

\bibitem{Rolfsen1976}
D.~Rolfsen,
\textit{Knots and Links},
Publish or Perish, Berkeley, 1976.

\end{thebibliography}

\end{document}
